% BuzzerModule
\section{【基本】BuzzerModule — deinit()パターン}

\begin{patternbox}{BuzzerModule で学べるパターン}
\begin{itemize}
    \item deinit()によるPWMリソースの明示的解放
    \item リクエスト方式による出力制御（requestPlayフラグ）
    \item ModuleTimerによる時間指定再生の自動停止
\end{itemize}
\end{patternbox}

\subsection{Data構造体のリクエストパターン}

\begin{lstlisting}[style=cppstyle, caption={BuzzerModule.h — Data}]
struct BuzzerData {
    uint16_t frequency = 0;       // 再生周波数 [Hz]（0=消音）
    unsigned long durationMs = 0; // 再生時間 [ms]（0=無制限）
    bool     requestPlay = false; // ロジックフェーズからの再生リクエスト
};
\end{lstlisting}

\begin{pointbox}[リクエストフラグパターン]
\texttt{requestPlay}は「ロジックフェーズからの指示を1回だけ処理する」ためのフラグ。
\texttt{updateOutput()}が処理後に\texttt{false}に戻すため、次ループでは再度処理しない。
LedModuleのように毎ループ同じ状態を出力するのではなく、
「イベント的に1回だけ何かをする」場合にこのパターンを使う。
\end{pointbox}

\subsection{updateOutput()の実装}

\begin{lstlisting}[style=cppstyle, caption={BuzzerModule.cpp — updateOutput()}]
void BuzzerModule::updateOutput(SystemData& data) {
    // 再生リクエスト処理
    if (data.buzzer.requestPlay) {
        data.buzzer.requestPlay = false;  // 処理済みフラグクリア

        if (data.buzzer.frequency > 0) {
            _toneOn(data.buzzer.frequency);
            if (data.buzzer.durationMs > 0) {
                _durationTimer.setTime();
                _timed = true;
            } else {
                _timed = false;  // 無制限再生
            }
        } else {
            _toneOff();
        }
    }

    // 時間指定ありの場合、経過したら自動消音
    if (_playing && _timed) {
        if (_durationTimer.getNowTime() >= data.buzzer.durationMs) {
            _toneOff();
        }
    }
}
\end{lstlisting}

\subsection{deinit()の実装}

\begin{lstlisting}[style=cppstyle, caption={BuzzerModule.cpp — deinit()}]
void BuzzerModule::deinit() {
    _toneOff();
    ledcDetachPin(_config.pin);  // PWMチャネルとピンの紐付けを解除
    Serial.println("[Buzzer] deinit OK");
}
\end{lstlisting}

\begin{pointbox}[deinit()の実装が必要なケース]
\texttt{IModule::deinit()}はデフォルトで空実装。
BuzzerModuleのように\textbf{ハードウェアリソース（PWMチャネル）を占有する}モジュールでは、
スリープ前やモジュール停止時にリソースを明示的に解放する必要がある。\\
\texttt{ledcDetachPin()}を呼ばないと、PWMチャネルが占有されたままスリープに入り、
復帰後に同じチャネルを使えなくなる可能性がある。
\end{pointbox}

\begin{notebox}[deinit()が不要なモジュール]
LedModuleやButtonModuleは\texttt{pinMode()}で設定したGPIOを使うだけで、
特別なリソースを占有しない。こうしたモジュールでは\texttt{deinit()}をオーバーライドしなくてよい。
\end{notebox}
